% 增量测度的三元分解：[0,1]上F(x)=x/2+(1/4)1_[1/2,1](x)+C(x)/4，C为Cantor函数。F(0)=0，F(1/2-)=3/8，F(1/2)=5/8，F(1)=1；F'=1/2几乎处处，导数积分1/2；跳跃1/4对应点质量，Cantor部分增量1/4。曲线坐标由同目录sample_cantor.py在步长1/243的网格上计算。横向1单位=7cm，纵向1单位=4cm。
\begin{tikzpicture}[x=7cm,y=4cm,
  lab/.style={font=\figurefont\fontsize{8}{10}\selectfont},
  note/.style={font=\FigureNoteFont\fontsize{7.5}{9}\selectfont,text=FigureMuted}]
\draw[FigureGrid,line width=.4pt] (0,0.25)--(1,0.25) (0,0.5)--(1,0.5) (0,0.75)--(1,0.75) (0,1)--(1,1);
\draw[FigureGreen,dashed,line width=.8pt] (0,0)--(1,0.5);
\draw[FigureBlue,line width=1.1pt] plot coordinates {(0.0000,0.0000) (0.0041,0.0099) (0.0082,0.0119) (0.0123,0.0218) (0.0165,0.0239) (0.0206,0.0259) (0.0247,0.0280) (0.0288,0.0378) (0.0329,0.0399) (0.0370,0.0498) (0.0412,0.0518) (0.0453,0.0539) (0.0494,0.0559) (0.0535,0.0580) (0.0576,0.0601) (0.0617,0.0621) (0.0658,0.0642) (0.0700,0.0662) (0.0741,0.0683) (0.0782,0.0782) (0.0823,0.0802) (0.0864,0.0901) (0.0905,0.0921) (0.0947,0.0942) (0.0988,0.0963) (0.1029,0.1061) (0.1070,0.1082) (0.1111,0.1181) (0.1152,0.1201) (0.1193,0.1222) (0.1235,0.1242) (0.1276,0.1263) (0.1317,0.1283) (0.1358,0.1304) (0.1399,0.1325) (0.1440,0.1345) (0.1481,0.1366) (0.1523,0.1386) (0.1564,0.1407) (0.1605,0.1427) (0.1646,0.1448) (0.1687,0.1469) (0.1728,0.1489) (0.1770,0.1510) (0.1811,0.1530) (0.1852,0.1551) (0.1893,0.1572) (0.1934,0.1592) (0.1975,0.1613) (0.2016,0.1633) (0.2058,0.1654) (0.2099,0.1674) (0.2140,0.1695) (0.2181,0.1716) (0.2222,0.1736) (0.2263,0.1835) (0.2305,0.1855) (0.2346,0.1954) (0.2387,0.1975) (0.2428,0.1995) (0.2469,0.2016) (0.2510,0.2115) (0.2551,0.2135) (0.2593,0.2234) (0.2634,0.2254) (0.2675,0.2275) (0.2716,0.2296) (0.2757,0.2316) (0.2798,0.2337) (0.2840,0.2357) (0.2881,0.2378) (0.2922,0.2398) (0.2963,0.2419) (0.3004,0.2518) (0.3045,0.2538) (0.3086,0.2637) (0.3128,0.2658) (0.3169,0.2678) (0.3210,0.2699) (0.3251,0.2797) (0.3292,0.2818) (0.3333,0.2917) (0.3374,0.2937) (0.3416,0.2958) (0.3457,0.2978) (0.3498,0.2999) (0.3539,0.3020) (0.3580,0.3040) (0.3621,0.3061) (0.3663,0.3081) (0.3704,0.3102) (0.3745,0.3122) (0.3786,0.3143) (0.3827,0.3164) (0.3868,0.3184) (0.3909,0.3205) (0.3951,0.3225) (0.3992,0.3246) (0.4033,0.3266) (0.4074,0.3287) (0.4115,0.3308) (0.4156,0.3328) (0.4198,0.3349) (0.4239,0.3369) (0.4280,0.3390) (0.4321,0.3410) (0.4362,0.3431) (0.4403,0.3452) (0.4444,0.3472) (0.4486,0.3493) (0.4527,0.3513) (0.4568,0.3534) (0.4609,0.3555) (0.4650,0.3575) (0.4691,0.3596) (0.4733,0.3616) (0.4774,0.3637) (0.4815,0.3657) (0.4856,0.3678) (0.4897,0.3699) (0.4938,0.3719) (0.4979,0.3740) (0.5000,0.3750)};
\draw[FigureBlue,line width=1.1pt] plot coordinates {(0.5000,0.6250) (0.5021,0.6260) (0.5062,0.6281) (0.5103,0.6301) (0.5144,0.6322) (0.5185,0.6343) (0.5226,0.6363) (0.5267,0.6384) (0.5309,0.6404) (0.5350,0.6425) (0.5391,0.6445) (0.5432,0.6466) (0.5473,0.6487) (0.5514,0.6507) (0.5556,0.6528) (0.5597,0.6548) (0.5638,0.6569) (0.5679,0.6590) (0.5720,0.6610) (0.5761,0.6631) (0.5802,0.6651) (0.5844,0.6672) (0.5885,0.6692) (0.5926,0.6713) (0.5967,0.6734) (0.6008,0.6754) (0.6049,0.6775) (0.6091,0.6795) (0.6132,0.6816) (0.6173,0.6836) (0.6214,0.6857) (0.6255,0.6878) (0.6296,0.6898) (0.6337,0.6919) (0.6379,0.6939) (0.6420,0.6960) (0.6461,0.6980) (0.6502,0.7001) (0.6543,0.7022) (0.6584,0.7042) (0.6626,0.7063) (0.6667,0.7083) (0.6708,0.7182) (0.6749,0.7203) (0.6790,0.7301) (0.6831,0.7322) (0.6872,0.7342) (0.6914,0.7363) (0.6955,0.7462) (0.6996,0.7482) (0.7037,0.7581) (0.7078,0.7602) (0.7119,0.7622) (0.7160,0.7643) (0.7202,0.7663) (0.7243,0.7684) (0.7284,0.7704) (0.7325,0.7725) (0.7366,0.7746) (0.7407,0.7766) (0.7449,0.7865) (0.7490,0.7885) (0.7531,0.7984) (0.7572,0.8005) (0.7613,0.8025) (0.7654,0.8046) (0.7695,0.8145) (0.7737,0.8165) (0.7778,0.8264) (0.7819,0.8284) (0.7860,0.8305) (0.7901,0.8326) (0.7942,0.8346) (0.7984,0.8367) (0.8025,0.8387) (0.8066,0.8408) (0.8107,0.8428) (0.8148,0.8449) (0.8189,0.8470) (0.8230,0.8490) (0.8272,0.8511) (0.8313,0.8531) (0.8354,0.8552) (0.8395,0.8573) (0.8436,0.8593) (0.8477,0.8614) (0.8519,0.8634) (0.8560,0.8655) (0.8601,0.8675) (0.8642,0.8696) (0.8683,0.8717) (0.8724,0.8737) (0.8765,0.8758) (0.8807,0.8778) (0.8848,0.8799) (0.8889,0.8819) (0.8930,0.8918) (0.8971,0.8939) (0.9012,0.9037) (0.9053,0.9058) (0.9095,0.9079) (0.9136,0.9099) (0.9177,0.9198) (0.9218,0.9218) (0.9259,0.9317) (0.9300,0.9338) (0.9342,0.9358) (0.9383,0.9379) (0.9424,0.9399) (0.9465,0.9420) (0.9506,0.9441) (0.9547,0.9461) (0.9588,0.9482) (0.9630,0.9502) (0.9671,0.9601) (0.9712,0.9622) (0.9753,0.9720) (0.9794,0.9741) (0.9835,0.9761) (0.9877,0.9782) (0.9918,0.9881) (0.9959,0.9901) (1.0000,1.0000)};
\fill[white] (0.5,0.375) circle (1.8pt); \draw[FigureBlue,line width=.7pt] (0.5,0.375) circle (1.8pt);
\fill[FigureBlue] (0.5,0.625) circle (1.8pt);
\draw[FigureRed,line width=.9pt,{Stealth[length=4pt]}-{Stealth[length=4pt]}] (0.53,0.385)--(0.53,0.615);
\node[note,anchor=west,text=FigureRed] at (0.545,0.5) {跳跃$\frac14$：点质量};
\draw[->,FigureStroke,line width=.7pt] (-0.02,0)--(1.07,0) node[lab,right] {$x$};
\draw[->,FigureStroke,line width=.7pt] (0,-0.02)--(0,1.12) node[lab,above] {$F(x)$};
\foreach \x/\t in {0/0,0.5/\frac12,1/1}{\draw[FigureStroke] (\x,0)--(\x,-0.015); \node[lab,below] at (\x,-0.01) {$\t$};}
\foreach \y/\t in {0.375/\frac38,0.625/\frac58,1/1}{\draw[FigureStroke] (0,\y)--(-0.01,\y); \node[lab,left] at (-0.008,\y) {$\t$};}
\node[note,anchor=west,text=FigureGreen] at (0.62,0.22) {虚线：绝对连续部分$x/2$};
\node[note,anchor=west] at (0.05,0.86) {细小台阶：Cantor部分共增加$\frac14$};
\end{tikzpicture}
